\section*{NeurIPS Paper Checklist}

\begin{enumerate}

\item {\bf Claims}
    \item[] Question: Do the main claims made in the abstract and introduction accurately reflect the paper's contributions and scope?
    \item[] Answer: \answerYes{}
    \item[] Justification: The abstract and Section~\ref{sec:system}--\ref{sec:visualizations} precisely describe the three technical contributions (aligned UMAP, velocity field, divergence) and their scope (visualisation tool for simple Gaussian ES). No performance claims beyond qualitative demonstration are made.

\item {\bf Limitations}
    \item[] Question: Does the paper discuss the limitations of the work performed by the authors?
    \item[] Answer: \answerYes{}
    \item[] Justification: Section~\ref{sec:limits} dedicates a full section to limitations, covering UMAP stochasticity, the non-Procrustes nature of the alignment, high-dimensional interpretation caveats, the restriction to simple ES, and the fidelity gap of the BRAX proxy.

\item {\bf Theory assumptions and proofs}
    \item[] Question: For each theoretical result, does the paper provide the full set of assumptions and a complete (and correct) proof?
    \item[] Answer: \answerNA{}
    \item[] Justification: The paper makes no theoretical claims; all contributions are algorithmic and empirical. Equations~\eqref{eq:align}--\eqref{eq:divergence} are definitions, not theorems.

\item {\bf Experimental result reproducibility}
    \item[] Question: Does the paper fully disclose all the information needed to reproduce the main experimental results of the paper to the extent that it affects the main claims and/or conclusions of the paper (regardless of whether the code and data are provided or not)?
    \item[] Answer: \answerYes{}
    \item[] Justification: Appendix~\ref{app:impl} provides pseudocode for the alignment algorithm, full dataset specifications, MLP architectures, ES hyperparameters, and the BRAX proxy dynamics equations. All random seeds are exposed as configurable parameters in the Streamlit UI.

\item {\bf Open access to data and code}
    \item[] Question: Does the paper provide open access to the data and code, with sufficient instructions to faithfully reproduce the main experimental results, as described in supplemental material?
    \item[] Answer: \answerYes{}
    \item[] Justification: The full codebase is released under the MIT licence with a \texttt{pixi}-managed reproducible environment. The command \texttt{pixi run app} reproduces all interactive figures; \texttt{pixi run figs} reproduces static publication figures. All datasets used are either generated programmatically (Make Moons), bundled with scikit-learn (Digits), implemented in pure NumPy (BRAX proxy), or downloaded from a stable public URL (CIFAR-10).

\item {\bf Experimental setting/details}
    \item[] Question: Does the paper specify all the training and test details (e.g., data splits, hyperparameters, how they were chosen, type of optimizer) necessary to understand the results?
    \item[] Answer: \answerYes{}
    \item[] Justification: Table~\ref{tab:datasets} and Appendix~\ref{app:impl} specify all hyperparameters ($G$, $N$, $\sigma$, $h$, $\lambda_{\text{align}}$, grid resolution). No gradient-based training occurs; the optimiser is the truncation-selection Gaussian ES defined in Equation~\eqref{eq:es}.

\item {\bf Experiment statistical significance}
    \item[] Question: Does the paper report error bars suitably and correctly defined or other appropriate information about the statistical significance of the experiments?
    \item[] Answer: \answerYes{}
    \item[] Justification: Table~\ref{tab:compute} reports mean $\pm$ standard deviation of wall-clock time over three independent seeds. The qualitative visualisation claims in Section~\ref{sec:experiments} are deliberately not quantified with significance tests, as the contribution is a visual diagnostic tool; multiple seed comparisons are discussed qualitatively.

\item {\bf Experiments compute resources}
    \item[] Question: For each experiment, does the paper provide sufficient information on the computer resources (type of compute workers, memory, time of execution) needed to reproduce the experiments?
    \item[] Answer: \answerYes{}
    \item[] Justification: Section~\ref{sec:experiments} states the hardware (Apple M1, 16\,GB RAM, no GPU) and Table~\ref{tab:compute} reports per-dataset wall-clock times at default parameters.

\item {\bf Code of ethics}
    \item[] Question: Does the research conducted in the paper conform, in every respect, with the NeurIPS Code of Ethics \url{https://neurips.cc/public/EthicsGuidelines}?
    \item[] Answer: \answerYes{}
    \item[] Justification: The work is a visualisation tool for scientific transparency. It uses only public benchmark datasets (Make Moons, Digits, CIFAR-10) and a synthetic proxy task. No human subjects, personal data, or sensitive content is involved.

\item {\bf Broader impacts}
    \item[] Question: Does the paper discuss both potential positive societal impacts and negative societal impacts of the work performed?
    \item[] Answer: \answerNA{}
    \item[] Justification: The work is a scientific diagnostic and visualisation tool for the neuroevolution research community. It carries no direct path to negative societal applications; its sole purpose is to increase transparency and interpretability of existing evolutionary algorithms.

\item {\bf Safeguards}
    \item[] Question: Does the paper describe safeguards that have been put in place for responsible release of data or models that have a high risk for misuse (e.g., pre-trained language models, image generators, or scraped datasets)?
    \item[] Answer: \answerNA{}
    \item[] Justification: The tool trains small classification and control MLPs on standard academic benchmarks. No large pre-trained models, scraped data, or generative components are released.

\item {\bf Licenses for existing assets}
    \item[] Question: Are the creators or original owners of assets (e.g., code, data, models), used in the paper, properly credited and are the license and terms of use explicitly mentioned and properly respected?
    \item[] Answer: \answerYes{}
    \item[] Justification: All libraries (NumPy, SciPy, scikit-learn, Matplotlib, Plotly, Streamlit, UMAP-learn) are cited in the references or text. CIFAR-10 is used in accordance with its non-commercial research licence. Make Moons and Digits are part of scikit-learn (BSD 3-Clause). The BRAX proxy is an original NumPy implementation inspired by, but not derived from, the JAX BRAX codebase.

\item {\bf New assets}
    \item[] Question: Are new assets introduced in the paper well documented and is the documentation provided alongside the assets?
    \item[] Answer: \answerYes{}
    \item[] Justification: The NeuroEvo-Viz codebase constitutes the primary new asset. Each module contains NumPy-style docstrings documenting parameters, return values, and design rationale. A \texttt{README.md} describes installation, available tasks, and usage commands.

\item {\bf Crowdsourcing and research with human subjects}
    \item[] Question: For crowdsourcing experiments and research with human subjects, does the paper include the full text of instructions given to participants and screenshots, if applicable, as well as details about compensation (if any)?
    \item[] Answer: \answerNA{}
    \item[] Justification: No crowdsourcing or human subjects research was conducted.

\item {\bf Institutional review board (IRB) approvals or equivalent for research with human subjects}
    \item[] Question: Does the paper describe potential risks incurred by study participants, whether such risks were disclosed to the subjects, and whether Institutional Review Board (IRB) approvals (or an equivalent approval/review based on the requirements of your country or institution) were obtained?
    \item[] Answer: \answerNA{}
    \item[] Justification: No human subjects are involved in this research.

\item {\bf Declaration of LLM usage}
    \item[] Question: Does the paper describe the usage of LLMs if it is an important, original, or non-standard component of the core methods in this research?
    \item[] Answer: \answerNA{}
    \item[] Justification: LLMs were not used as a component of the core methodology. They were used solely for writing assistance and do not affect the scientific contributions or experimental results.

\end{enumerate}
